%% Advances in Combinatorics article template
%%
%% aic-template.tex v0.33
%%
%% AUTHOR: Fill in fields (or see warnings) below marked with "AUTHOR"
%% ** Add as few macro / package definitions as possible
%% ** Compile with "pdflatex"; make sure that
%%           aic.cls and tocbase.cls are in the same directory.
%%
%% EDITOR: Fill in fields below marked with "EDITOR"
%%    and check that authors properly filled in field marked with "AUTHOR"


\documentclass{aic}

%%%%%%%%%%%%%%%%%%%% Packages and macros
\usepackage{amsmath,amsthm,amssymb,hyperref, color, bm}
\usepackage{blkarray}

\newcommand{\remove}[1]{}

\makeatletter
\newtheorem*{rep@theorem}{\rep@title}
\newcommand{\newreptheorem}[2]{%
\newenvironment{rep#1}[1]{%
 \def\rep@title{#2 \ref{##1}}%
 \begin{rep@theorem}}%
 {\end{rep@theorem}}}
\makeatother

% \hypersetup{
% 	colorlinks=true,
% 	linkcolor=blue,%
% 	citecolor=blue,
% 	linktoc=page
% }
\newcommand\numberthis{\addtocounter{equation}{1}\tag{\theequation}} 
\renewcommand{\thefootnote}{\fnsymbol{footnote}}
  

\newtheorem{thm}{Theorem}[section]
\newreptheorem{thm}{Theorem}
\newtheorem{claim}[thm]{Claim}
\newtheorem{lem}[thm]{Lemma}
\newreptheorem{lem}{Lemma}
\newtheorem{define}[thm]{Definition}
\newtheorem{cor}[thm]{Corollary}
\newtheorem{obs}[thm]{Observation}
\newtheorem{stat}[thm]{Statement}
\newtheorem{example}[thm]{Example}
\newtheorem{comm}[thm]{Comment}
\newtheorem{construction}[thm]{Construction}
\newtheorem{conjecture}[thm]{Conjecture}
\newtheorem{THM}{Theorem}
\newtheorem{question}{Question}
\newtheorem{fact}[thm]{Fact}
\newtheorem{prop}[thm]{Proposition}




\def\dst{\displaystyle}

% Special commands

\def\L{{\mathcal L}}
\def\B{{\mathcal B}}
\def\GL{\text{GL}}
\def\F{{\mathbb{F}}}
\def\E{{\mathbb{E}}}
\def\kF{{\overline{\mathbb{F}}}}
\def\kS{{\overline{S}}}
\def\fS{{\mathfrak{S}}}
\def\kR{{\overline{R}}}
\def\ce{{\hat{e}}}
\def\Q{{\mathbb{Q}}}
\def\Z{{\mathbb{Z}}}
\def\N{{\mathbb{N}}}
\def\R{{\mathbb{R}}}
\def\K{{\mathbb{K}}}
\def\C{{\mathbb{C}}}
\def\A{{\mathbb{A}}}
\def\P{{\mathbb{P}}}
\def\cP{{\cal P}}
\def\cB{{\cal B}}
\def\cK{{\cal K}}
\def\cS{{\cal S}}
\def\cH{{\cal H}}
\def\cL{{\cal L}}
\def\cG{{\cal G}}
\def\cC{{\cal C}}
\def\cF{{\cal F}}
\def\cO{{\cal O}}
\def\cJ{{\cal J}}
\def\cU{{\cal U}}
\def\cM{\mathcal{M}}
\def\cR{\mathcal{R}}
\def\cS{{\mathcal S}}
\def\cE{{\mathcal E}}
\def\V{{\mathbf{V}}}
\def\I{{\mathbf{I}}}
\def\Ind{{\mathbf{1}}}
\def\bx{{\mathbf x}}
\def\bp{{\mathbf p}}
\def\bj{{\mathbf j}}
\def\balpha{{\bm \alpha}}
\def\bbeta{{\bm\beta}}
\def\bq{{\mathbf q}}
\def\by{{\mathbf y}}
\def\E{{\mathbb E}}


\def\half{ \frac{1}{2}}
\def\sumN{\sum_{i=1}^n}
\def\_{\,\,\,\,\,}
\def\prob{{\mathbf{Pr}}}
\def\modulo{\text{mod}}
\def\omm{ \{0,1\} }
\def\id{ \textit{id} }

\def\then{\Rightarrow}

\def\D{{\partial}}
\def\gap{\textsf{gap}}
\def\com{\textsf{com}}
\def\sign{\textsf{sign}}
\def\supp{\textsf{supp}}
 \def\spar{\textsf{sparse}}
\def\tr{\textsf{tr}}
\def\Part{\textbf{Part}}
\def\reach{\textsf{Rea}}
\def\Mon{\textbf{Mon}}
\def\sing{\textbf{sing}}
\def\Und{\textsf{Und}}
\def\Comp{\textsf{Comp}}
\def\rank{\textsf{rank}}
\def\crank{\textsf{crank}}
\def\poly{\textsf{poly}}
\def\codim{\textsf{codim}}
\def\dim{\textsf{dim}}
\def\cp{\textsf{cp}}
\def\uni{\textsf{Uni}}
\def\ext{\textsf{\bf Ext}}
\def\extt{\textsf{\bf Ext2}}
\def\Cor{\textsf{Cor}}
\def\capac{\textsf{cap}}
\def\row{\textsf{Row}}
\def\Dir{\text{Dir}}
\def\Coeff{\text{Coeff}}
\def\lt{\ltimes}
\def\inj{\hookrightarrow}
\def\es{\emptyset}



\def\vc{\textsf{VC-dim}}

\def\UG{{\cal G}}
\def\DG{{\cal D}}

\newcommand{\eps}{\epsilon}
% \newcommand\leftidx[3]{%
%   {\vphantom{#2}}#1#2#3%
% }
\usepackage{leftidx}

\newcommand{\red}[1]{\textcolor{red}{#1}}



%%%%%%%%%%%%TIKZ PACKAGE FOR FIGURES%%%%%%%%%%%%%%
\usepackage{tikz-cd}
% \usepackage{pgf,tikz,pgfplots}
% \pgfplotsset{compat=1.15}
% \usepackage{mathrsfs}
% \usetikzlibrary{arrows}
%%%%%%%%%%%%%%%%%%%%%%%%%%%%%%%%%%%%%%%%%%%%%%%%%%%
%%%%%%%%%%%%%%%%%%%%%%%%%%%%%%%%



%%%%%%%%%%%%%%%%%%%%%%%%%%%%%%%%%%%%%%%%%%%%%%%%
%% AUTHOR: Fill in meta-data below:
\aicAUTHORdetails{%
  title = {The Kakeya Set conjecture over $\mathbb{Z}/N\mathbb{Z}$ for general $N$}, %% please capitalize all significant words
  author = {Manik Dhar},
    %% Please use the format for commas as follows:
    %% "A", or "A and B", or "A, B, and C", or "A, B, C, and D", etc.
  plaintextauthor = {Manik Dhar},
    %% An author list in plain text: Use the format
    %% "A", or "A, B", or "A, B, C", etc.
    %% NOTE: No LaTeX code in author names.
    %% NOTE: No "and" at the end--simply comma separated,
    % 
 %% The remaining lines in this section are optional:
    %
    %% IF YOUR TITLE CONTAINS MATH OR LATEX such as accented characters: 
    %% Add a "plain text title";  otherwise comment out the next line:
  plaintexttitle = {The Kakeya Set conjecture over Z/N Z for general N}, %%  title without math or LaTeX
    %
    %% ONLY IF YOUR TITLE IS TOO LONG to fit in the page headers, please 
    %% add an abbreviated version of the title, otherwise comment it out:
  % runningtitle = {R\"odl's $n^{\log\log n}$ Bound}, 
    %
    %% ONLY IF YOUR AUTHOR LIST IS TOO LONG to fit in the page headers, 
    %% add an abbreviated version, otherwise comment it out:
  % runningauthor = {Paul Erd\H{o}s, Johan H{\aa}stad, L\'aszl\'o Lov\'asz, and Andrew C-C. Yao},
    %% you can replace first names and/or middle names with initials.
    %
    %% ONLY IF YOUR AUTHOR LIST IS TOO LONG to fit the copyright entry
    %% on the bottom of the front page,
    %% add an abbreviated version, otherwise comment it out:
  % copyrightauthor = {Manik Dhar},
    %% Note that the copyrightauthor  field will seldom be necessary;
    %% for instance, in this example with four authors, it would be 
    %% all right to comment it out and have all authors' full names 
    %% appear on the Copyright line
   %
   %% Include keywords of your choice: comma separated, lower case;
   %% comment out the "keywords" line if you don't wish to provide them
  keywords = {Kakeya conjecture, polynomial method},
}   %%% END \aicAUTHORdetails

%%%%%%%%%%%%%%%%%%%%%%%%%%%%%%%%%%%%%%%%%%%%%%%%
%%% EDITOR: please fill in the following data:
\aicEDITORdetails{%
   year={2024},
 %  volume={XX},
   number={2},
   received={18 January 2022},   % received date: example: 7 January 2017
   revised={4 June 2023},    % Optional revised date (you may comment it out)
   published={26 January 2024},  % published date
   doi={10.19086/aic.2024.2},      % XXX = number of paper, e.g. aic006 for paper#6
%                              % or  aic0006 (length of string arbitrary)
}   %%% END \aicEDITORdetails

\begin{document}

\begin{frontmatter}[classification=text]
%% EDITOR: this will force the keywords to appear right after the Abstract.
%%   If the abstract is too long and would force the keywords off the
%%   front page, please comment out % [classification=text] above
%%   This way the keywords will be floated on the bottom of the first page
%%   even though the Abstract spills over to the next page.

%%% AUTHOR: Title goes here.  This line is optional.  You must use it
%%   if title has footnote attached or requires nontrivial typesetting,
%%   e.g., inclusion of linebreaks to force nice layout.
\title{The Kakeya Set conjecture over $\mathbb{Z}/N\mathbb{Z}$ for general $N$} %% please capitalize all significant words

%%% AUTHOR:
%%% List all authors. If you wish, place grant acknowledgements in \thanks.
%%% In brackets include a short tag for each author.
\author[md]{Manik Dhar}
% \thanks{Department of Mathematics, Massachusetts institute of Technology. Email: \texttt{dmanik@mit.edu}. }


%%% AUTHOR: Abstract goes here
\begin{abstract}
We prove the Kakeya set conjecture for $\mathbb{Z}/N\mathbb{Z}$ for general $N$ as stated by Hickman and Wright~\cite{hickman2018fourier}. This entails extending and combining the techniques of Arsovski~\cite{arsovski2021padic} for $N=p^k$ and the author and Dvir~\cite{dhar2021proof} for the case of square-free $N$. We also prove stronger lower bounds for the size of $(m,\epsilon)$-Kakeya sets over $\Z/p^k\Z$ by extending the techniques of \cite{arsovski2021padic} using multiplicities as was done in \cite{saraf2008,dvir2013extensions}. In addition, we show our bounds are almost sharp by providing a new construction for Kakeya sets over $\Z/p^k\Z$ and $\Z/N\Z$.
\end{abstract}
\end{frontmatter}

%%% AUTHOR: body of paper starts here
\input{body.tex}

%%% AUTHOR: optional appendix here
\input{appendix}

%%% AUTHOR: optional acknowledgments here
% \section*{Acknowledgments} %%  you may comment this out if no Ackno
% The authors are grateful to the anonymous reviewers for finding
% a bug in the main result.

%%% AUTHOR:
%%% Bibliography goes here. Note that the arXiv cannot process bibtex
%%% or biber bibliographies.  Example of acceptable bibliograpy format:
\bibliographystyle{amsplain}
\begin{thebibliography}{10}
\providecommand{\bysame}{\leavevmode\hbox to3em{\hrulefill}\thinspace}
\providecommand{\MR}{\relax\ifhmode\unskip\space\fi MR }
% \MRhref is called by the amsart/book/proc definition of \MR.
\providecommand{\MRhref}[2]{%
  \href{http://www.ams.org/mathscinet-getitem?mr=#1}{#2}
}
\providecommand{\href}[2]{#2}
\bibitem{arsovski2021padic}
Bodan Arsovski, \emph{The p-adic {K}akeya conjecture}, arXiv preprint
  2108.03750v1, 08-08-2021.

\bibitem{arsovskiNew}
\bysame, \emph{The p-adic {K}akeya conjecture}, Journal of the American Mathematical Society, \textbf{37} (2023), 69-80.



\bibitem{bukh2021sharp}
Boris Bukh and Ting-Wei Chao, \emph{Sharp density bounds on the finite field
  {K}akeya}, Discrete Analysis 2021, no. 26.

\bibitem{CML_2018__10_1_3_0}
Xavier Caruso, \emph{Almost all non-archimedean {K}akeya sets have measure
  zero}, Confluentes Mathematici \textbf{10} (2018), no.~1, 3--40 (en).

\bibitem{dhar2022maximal}
Manik Dhar, \emph{{Maximal Kakeya and $(m,\epsilon)$-Kakeya bounds over
  $\mathbb{Z}/\text{{N}}\mathbb{Z}$ for general N}}, preprint arXiv:2209.11443
  (2022).

\bibitem{dhar2021proof}
Manik Dhar and Zeev Dvir, \emph{Proof of the {K}akeya set conjecture over rings
  of integers modulo square-free {N}}, Combinatorial Theory \textbf{1} (2021).

\bibitem{dhar2019simple}
Manik Dhar, Zeev Dvir, and Ben Lund, \emph{Simple proofs for furstenberg sets
  over finite fields}, Discrete Analysis (2021), no.~22 (en).

\bibitem{dummithablicsek2013}
Evan~P. Dummit and Márton Hablicsek, \emph{Kakeya sets over non-archimedean
  local rings}, Mathematika \textbf{59} (2013), no.~2, 257–266.

\bibitem{dvir2009size}
Zeev Dvir, \emph{On the size of {K}akeya sets in finite fields}, Journal of the
  American Mathematical Society \textbf{22} (2009), no.~4, 1093--1097.

\bibitem{dvir2013extensions}
Zeev Dvir, Swastik Kopparty, Shubhangi Saraf, and Madhu Sudan, \emph{Extensions
  to the method of multiplicities, with applications to {K}akeya sets and
  mergers}, SIAM Journal on Computing \textbf{42} (2013), no.~6, 2305--2328.

\bibitem{ellenberg2010kakeya}
Jordan~S Ellenberg, Richard Oberlin, and Terence Tao, \emph{The {K}akeya set
  and maximal conjectures for algebraic varieties over finite fields},
  Mathematika \textbf{56} (2010), no.~1, 1--25.

\bibitem{Fraser_2016}
Robert Fraser, \emph{Kakeya-type sets in local fields with finite residue
  field}, Mathematika \textbf{62} (2016), no.~2, 614--629.

\bibitem{zbMATH02610444}
G.~H. {Hardy} and S.~{Ramanujan}, \emph{{The normal number of prime factors of
  a number \(n\)}}, {Quart. J.} \textbf{48} (1917), 76--92 (English).

\bibitem{Hardy2008AnIT}
G.H. {Hardy} and E.M. Wright, \emph{An introduction to the theory of numbers},
  4th ed., 1975.

\bibitem{hickman2018fourier}
Jonathan Hickman and James Wright, \emph{The {F}ourier restriction and {K}akeya
  problems over rings of integers modulo {N}}, Discrete Analysis (2018),
  no.~11.

\bibitem{Lucas}
Edouard Lucas, \emph{Théorie des fonctions numériques simplement
  périodiques}, American Journal of Mathematics \textbf{1} (1878), no.~4,
  289--321.

\bibitem{saraf2008}
Shubhangi Saraf and Madhu Sudan, \emph{An improved lower bound on the size of
  {K}akeya sets over finite fields}, Anal. PDE \textbf{1} (2008), no.~3,
  375--379.

\bibitem{wolf1999}
Thomas Wolff, \emph{Recent work connected with the {K}akeya problem}, Prospects
  in mathematics (Princeton,NJ, 1996) (1999), 29--162.

\end{thebibliography}

%\bibliography{bibliography}
% \begin{thebibliography}{99}
% \bibitem{bergelson-johnson-moreira}
% Vitaly Bergelson, John H. Johnson Jr., and Joel Moreira.
% \newblock New polynomial and multidimensional extensions of classical partition
%   results.
% \newblock 2015, arXiv:1501.02408.

% \bibitem{cilleruelo}
% Javier Cilleruelo.
% \newblock Combinatorial problems in finite fields and {S}idon sets.
% \newblock {\em Combinatorica}, 32(5):497--511, 2012.

% \end{thebibliography}
%% AUTHOR: You can generate such a bibliography from a .bib file by 
%% running pdflatex/bibtex/pdflatex/pdflatex and then pasting the .bbl file
%% between \begin{thebibliography} and \end{bibliography}

%%% AUTHOR: Include a short description of each author following the
%%% structure below. Use the same short tags used previously.  
%%% Use \imageat{} and \imagedot{} instead of "@" and "." in
%%% email addresses-this replaces the symbols with graphics to avoid 
%%% e-mail address harvesting from the .pdf file
\begin{aicauthors}
\begin{authorinfo}[pgom]
  Manik Dhar\\
  Massachusetts Institute of Technology\\
  
  Cambridge, MA, USA\\
  dmanik\imageat{}mit\imagedot{}edu \\
  \url{https://dharmanik.github.io}
\end{authorinfo}
\end{aicauthors}

\end{document}

%%% Local Variables:
%%% mode: latex
%%% TeX-master: "aic-template"
%%% End:
